\documentclass[12pt,a4paper]{article}
\usepackage[utf8]{inputenc}
\usepackage[T1]{fontenc}
\usepackage[french]{babel}
\usepackage{geometry}
\usepackage{enumitem}
\geometry{margin=2.5cm}

\title{Compte rendu des modifications\\
\large Alignement avec le cahier des charges}
\author{Projet : Gestion de stock supermarché}
\date{\today}

\begin{document}

\maketitle

\section*{1. Contexte}
Dans la version initiale du cahier des charges, les fonctionnalités principales étaient bien identifiées (recherche, filtres/tri, journalisation), mais leur emplacement et leur comportement dans l'interface n'étaient pas suffisamment précisés.

Les dernières modifications ont donc eu pour objectif de lever ces ambiguïtés UX tout en respectant les besoins fonctionnels attendus.

\section*{2. Modifications effectuées}

\subsection*{2.1 Recherche de produit}
\begin{itemize}[leftmargin=1.2cm]
  \item Intégration d'une recherche accessible par \textbf{bouton}.
  \item Ouverture de la recherche dans une \textbf{fenêtre dédiée (modal)}.
  \item Affichage des \textbf{résultats en temps réel} pendant la saisie.
  \item Déploiement de cette logique sur l'accueil et la page liste produits pour homogénéiser l'expérience utilisateur.
\end{itemize}

\subsection*{2.2 Filtres et tri}
\begin{itemize}[leftmargin=1.2cm]
  \item Mise en place du filtrage avancé : catégorie, prix min/max, quantité min/max, disponibilité.
  \item Ajout du tri : nom, prix, quantité (ordre croissant/décroissant).
  \item Regroupement des filtres dans une \textbf{modal dédiée ``Filtres \& tri''} au lieu d'un affichage permanent.
  \item Repositionnement du bouton ``Filtres \& tri'' \textbf{juste avant le tableau des produits} (aligné à gauche), conformément à la clarification demandée.
\end{itemize}

\subsection*{2.3 Export CSV}
\begin{itemize}[leftmargin=1.2cm]
  \item Ajout d'un bouton \textbf{Export CSV} dans la page produits.
  \item Implémentation backend d'un endpoint d'export CSV.
  \item Export basé sur les filtres et tris actifs pour produire un fichier cohérent avec la vue utilisateur.
\end{itemize}

\subsection*{2.4 Journalisation}
\begin{itemize}[leftmargin=1.2cm]
  \item Création d'un mécanisme de journalisation des actions stock.
  \item Enregistrement des actions \textbf{CREATE / UPDATE / DELETE} côté backend.
  \item Ajout d'une table de logs en base pour l'historique.
  \item Affichage du journal dans l'interface pour consultation rapide.
\end{itemize}

\subsection*{2.5 Responsive et ergonomie}
\begin{itemize}[leftmargin=1.2cm]
  \item Réorganisation responsive des boutons d'action (recherche, export, ajout) en disposition verticale selon la zone d'interface.
  \item Simplification visuelle de la page liste en sortant les paramètres avancés de la vue principale.
\end{itemize}

\section*{3. Justification par rapport au cahier des charges}
Ces modifications ne changent pas le périmètre fonctionnel demandé ; elles précisent \textbf{l'implantation} des fonctionnalités dans l'interface. Elles répondent à un besoin de clarification du cahier des charges sur le \emph{où} et le \emph{comment}.

En résumé :
\begin{itemize}[leftmargin=1.2cm]
  \item Les fonctionnalités demandées sont bien présentes.
  \item Leur intégration UI a été ajustée pour améliorer la lisibilité, la cohérence et l'usage en contexte réel.
  \item Le comportement final est plus conforme à une utilisation opérationnelle (MVP soutenable en démonstration).
\end{itemize}

\section*{4. Conclusion}
Le projet est désormais aligné sur les exigences fonctionnelles du cahier des charges, avec une meilleure définition UX des éléments initialement ambiguës (recherche, filtres/tri, journalisation).

\end{document}